\documentclass{article}
\usepackage{listings}
\begin{document}
\section*{Задача: функция среднего арифметического}

Условие: написать функцию \texttt{mean(values)}, которая возвращает среднее
арифметическое списка чисел. Если список пустой, функция должна возвращать
\texttt{None}.

\begin{lstlisting}[language=Python]
def mean(values: list[float]) -> float | None:
    if not values:
        return None
    return sum(values) / len(values)


# Тесты
assert mean([1, 2, 3, 4, 5]) == 3.0
assert mean([10]) == 10.0
assert mean([]) is None
print("OK")
\end{lstlisting}

\textbf{Идея:} проверяю пустоту, возвращаю None. Иначе делю сумму на длину.
\end{document}
